% Options for packages loaded elsewhere
\PassOptionsToPackage{unicode}{hyperref}
\PassOptionsToPackage{hyphens}{url}
\documentclass[
]{article}
\usepackage{xcolor}
\usepackage[margin=1in]{geometry}
\usepackage{amsmath,amssymb}
\setcounter{secnumdepth}{-\maxdimen} % remove section numbering
\usepackage{iftex}
\ifPDFTeX
  \usepackage[T1]{fontenc}
  \usepackage[utf8]{inputenc}
  \usepackage{textcomp} % provide euro and other symbols
\else % if luatex or xetex
  \usepackage{unicode-math} % this also loads fontspec
  \defaultfontfeatures{Scale=MatchLowercase}
  \defaultfontfeatures[\rmfamily]{Ligatures=TeX,Scale=1}
\fi
\usepackage{lmodern}
\ifPDFTeX\else
  % xetex/luatex font selection
\fi
% Use upquote if available, for straight quotes in verbatim environments
\IfFileExists{upquote.sty}{\usepackage{upquote}}{}
\IfFileExists{microtype.sty}{% use microtype if available
  \usepackage[]{microtype}
  \UseMicrotypeSet[protrusion]{basicmath} % disable protrusion for tt fonts
}{}
\makeatletter
\@ifundefined{KOMAClassName}{% if non-KOMA class
  \IfFileExists{parskip.sty}{%
    \usepackage{parskip}
  }{% else
    \setlength{\parindent}{0pt}
    \setlength{\parskip}{6pt plus 2pt minus 1pt}}
}{% if KOMA class
  \KOMAoptions{parskip=half}}
\makeatother
\usepackage{color}
\usepackage{fancyvrb}
\newcommand{\VerbBar}{|}
\newcommand{\VERB}{\Verb[commandchars=\\\{\}]}
\DefineVerbatimEnvironment{Highlighting}{Verbatim}{commandchars=\\\{\}}
% Add ',fontsize=\small' for more characters per line
\usepackage{framed}
\definecolor{shadecolor}{RGB}{248,248,248}
\newenvironment{Shaded}{\begin{snugshade}}{\end{snugshade}}
\newcommand{\AlertTok}[1]{\textcolor[rgb]{0.94,0.16,0.16}{#1}}
\newcommand{\AnnotationTok}[1]{\textcolor[rgb]{0.56,0.35,0.01}{\textbf{\textit{#1}}}}
\newcommand{\AttributeTok}[1]{\textcolor[rgb]{0.13,0.29,0.53}{#1}}
\newcommand{\BaseNTok}[1]{\textcolor[rgb]{0.00,0.00,0.81}{#1}}
\newcommand{\BuiltInTok}[1]{#1}
\newcommand{\CharTok}[1]{\textcolor[rgb]{0.31,0.60,0.02}{#1}}
\newcommand{\CommentTok}[1]{\textcolor[rgb]{0.56,0.35,0.01}{\textit{#1}}}
\newcommand{\CommentVarTok}[1]{\textcolor[rgb]{0.56,0.35,0.01}{\textbf{\textit{#1}}}}
\newcommand{\ConstantTok}[1]{\textcolor[rgb]{0.56,0.35,0.01}{#1}}
\newcommand{\ControlFlowTok}[1]{\textcolor[rgb]{0.13,0.29,0.53}{\textbf{#1}}}
\newcommand{\DataTypeTok}[1]{\textcolor[rgb]{0.13,0.29,0.53}{#1}}
\newcommand{\DecValTok}[1]{\textcolor[rgb]{0.00,0.00,0.81}{#1}}
\newcommand{\DocumentationTok}[1]{\textcolor[rgb]{0.56,0.35,0.01}{\textbf{\textit{#1}}}}
\newcommand{\ErrorTok}[1]{\textcolor[rgb]{0.64,0.00,0.00}{\textbf{#1}}}
\newcommand{\ExtensionTok}[1]{#1}
\newcommand{\FloatTok}[1]{\textcolor[rgb]{0.00,0.00,0.81}{#1}}
\newcommand{\FunctionTok}[1]{\textcolor[rgb]{0.13,0.29,0.53}{\textbf{#1}}}
\newcommand{\ImportTok}[1]{#1}
\newcommand{\InformationTok}[1]{\textcolor[rgb]{0.56,0.35,0.01}{\textbf{\textit{#1}}}}
\newcommand{\KeywordTok}[1]{\textcolor[rgb]{0.13,0.29,0.53}{\textbf{#1}}}
\newcommand{\NormalTok}[1]{#1}
\newcommand{\OperatorTok}[1]{\textcolor[rgb]{0.81,0.36,0.00}{\textbf{#1}}}
\newcommand{\OtherTok}[1]{\textcolor[rgb]{0.56,0.35,0.01}{#1}}
\newcommand{\PreprocessorTok}[1]{\textcolor[rgb]{0.56,0.35,0.01}{\textit{#1}}}
\newcommand{\RegionMarkerTok}[1]{#1}
\newcommand{\SpecialCharTok}[1]{\textcolor[rgb]{0.81,0.36,0.00}{\textbf{#1}}}
\newcommand{\SpecialStringTok}[1]{\textcolor[rgb]{0.31,0.60,0.02}{#1}}
\newcommand{\StringTok}[1]{\textcolor[rgb]{0.31,0.60,0.02}{#1}}
\newcommand{\VariableTok}[1]{\textcolor[rgb]{0.00,0.00,0.00}{#1}}
\newcommand{\VerbatimStringTok}[1]{\textcolor[rgb]{0.31,0.60,0.02}{#1}}
\newcommand{\WarningTok}[1]{\textcolor[rgb]{0.56,0.35,0.01}{\textbf{\textit{#1}}}}
\usepackage{graphicx}
\makeatletter
\newsavebox\pandoc@box
\newcommand*\pandocbounded[1]{% scales image to fit in text height/width
  \sbox\pandoc@box{#1}%
  \Gscale@div\@tempa{\textheight}{\dimexpr\ht\pandoc@box+\dp\pandoc@box\relax}%
  \Gscale@div\@tempb{\linewidth}{\wd\pandoc@box}%
  \ifdim\@tempb\p@<\@tempa\p@\let\@tempa\@tempb\fi% select the smaller of both
  \ifdim\@tempa\p@<\p@\scalebox{\@tempa}{\usebox\pandoc@box}%
  \else\usebox{\pandoc@box}%
  \fi%
}
% Set default figure placement to htbp
\def\fps@figure{htbp}
\makeatother
\setlength{\emergencystretch}{3em} % prevent overfull lines
\providecommand{\tightlist}{%
  \setlength{\itemsep}{0pt}\setlength{\parskip}{0pt}}
\usepackage{framed}
\usepackage{xcolor}
\let\oldquote\quote
\let\endoldquote\endquote
\renewenvironment{quote}{\begin{leftbar}}{\end{leftbar}}
\usepackage{bookmark}
\IfFileExists{xurl.sty}{\usepackage{xurl}}{} % add URL line breaks if available
\urlstyle{same}
\hypersetup{
  pdftitle={Review 2: Author responses},
  pdfauthor={Panis \& Ramsey},
  hidelinks,
  pdfcreator={LaTeX via pandoc}}

\title{Review 2: Author responses}
\author{Panis \& Ramsey}
\date{07.10.2025}

\begin{document}
\maketitle

\section{Editor}\label{editor}

\begin{quote}
AMPPS - Decision on Manuscript ID AMPPS-24-0232.R1 Advances in Methods
and Practices in Psychological Science 25-Sep-2025 Dear Dr.~Panis: Thank
you for re-submitting your Tutorial (AMPPS-24-0232.R1) entitled ``Event
History Analysis for psychological time-to-event data: A tutorial in R
with examples in Bayesian and frequentist workflows'' to Advances in
Methods and Practices in Psychological Science (AMPPS). I was fortunate
to receive feedback from three of the original reviewers, and you will
find their comments below this email and attached. As you will see, all
three reviewers believe your paper is improved and that this work
continues to have the potential to make a strong contribution. I
continue to share this assessment, as well as the need for some further
refinements. Therefore, I again invite you to revise and resubmit your
manuscript for further consideration at AMPPS. All three reviewers
acknowledge that the revised manuscript has improved in clarity,
structure, and focus, with Reviewer 1 especially commending its sharper
aims and increased utility for experimental psychologists working with
response time data.
\end{quote}

\begin{quote}
That said, accessibility remains a shared concern. Reviewer 3 highlights
that the tutorial is still challenging to follow, particularly for
readers less familiar with event history analysis (EHA) or
speed--accuracy tradeoff (SAT) approaches, and urges that key concepts
and research design features be defined and explained in the main text
rather than deferred to the supplement.
\end{quote}

Response: We define and explain key concepts and research design
features in the main text now.

\begin{quote}
Reviewer 2 raises several clarity-related concerns--particularly about
the bin width discussion and about the interpretability of modeling
decisions such as time window selection.
\end{quote}

Response: We include a new supplementary figure to illustrate the choice
of bin width.

\begin{quote}
Both Reviewers 2 and 3 emphasize the need to better motivate
methodological choices, especially the use of discrete-time EHA, which
is currently justified by lab familiarity rather than a clear
articulation of its analytic tradeoffs (R3, citing p.~8 of the ms).
\end{quote}

Response: We now articulate the analytic tradeoffs when motivating our
use of discrete-time EHA.

\begin{quote}
Reviewer 3 also recommends that the paper more clearly identify the
types of research questions EHA is particularly well-suited to address,
ideally with earlier and more focused examples (R3, citing p.~10 of the
ms). I agree with this point and feel it's critical to make the paper as
accessible as possible-- perhaps with (and not limited to) easier
examples up front, a glossary of sorts, visualizations; anything (and
everything) that will make this complex paper more inviting.
\end{quote}

Response: We include a glossary, visualizations (previous section A of
the supplementary material), and example research questions in the main
text of the revised ms.

\begin{quote}
While the paper is technically strong and substantively valuable, its
didactic function would be enhanced by more clearly integrating
definitions, rationales, and analytic guidance into the main text for
the broader AMPPS readership. I'd encourage you to go through the paper
start-to-end with an eye toward improving all possible points of entry
and overall accessibility.
\end{quote}

Response: We tried to improve all possible points of entry, as well as
overall accessibility.

\begin{quote}
If you choose to submit a revision, please include a letter detailing
your point-by-point responses to each reviewer comment and indicating
how you changed the manuscript to address them. The revised manuscript
may undergo further peer review. We ask that you submit your revision
within three months. Please let us know if you will not be able to meet
this deadline. To submit your revision, log into
\url{http://mc.manuscriptcentral.com/ampps} and enter your Author
Center, where you will find your manuscript title listed under
``Manuscripts with Decisions.'' Under ``Actions,'' click on ``Create a
Revision.'' Your manuscript number has been appended to denote a
revision. IMPORTANT: Your original files are available to you when you
upload your revised manuscript. Please delete any redundant or outdated
files before completing the submission. Once again, thank you for
submitting your manuscript to AMPPS and I look forward to receiving your
revision.
\end{quote}

\section{Reviewer: 1}\label{reviewer-1}

\begin{quote}
Comments to the Author From my perspective, the authors have been
extremely responsive to previous feedback. The manuscript has improved
considerably. The manuscript now has a much clearer focus in terms of
its aims (i.e., by targeting readers from experimental and cognitive
psychology working with response time data), which greatly increases the
utility of this work. The manuscript is also organized in a much more
logical way, and it does a great job of referring to and contextualizing
other helpful resources. This work has the potential to make a major
contribution to the field and it has been a pleasure to review this
work.
\end{quote}

Response: We thank this reviewer for their effort and comments.

\section{Reviewer 2}\label{reviewer-2}

\begin{quote}
Reviewer report for AMPPS-24-0232: ``Event History Analysis for
psychological time-to- event data: A tutorial in R with examples in
Bayesian and frequentist workflows'' Summary: I appreciate the effort
taken by the authors to strengthen the manuscript and find the revised
version much improved in terms of clarity and accessibility. In
particular, the introduction and motivation are much easier to follow. I
believe this manuscript is a useful contribution to the literature;
however, I have a few remaining comments.
\end{quote}

Response: Thank you for your revision, and your recognition of our
efforts.

\begin{quote}
Major comments:
\end{quote}

\begin{quote}
\begin{enumerate}
\def\labelenumi{\arabic{enumi}.}
\tightlist
\item
  I appreciate the authors' attempt to address previous comments on
  determining the appropriate bin width (p.~12). The authors note that
  bin widths need not be constant over time was quite useful. This new
  paragraph on bin width has been added under the section ``Implications
  for research design'' and I wonder if it would be more appropriately
  placed elsewhere. The text's current location implies that bin width
  should be determined when designing the experiment, rather than during
  analysis. I worry that suggesting that bin width be determined during
  the design of the experiment could suggest (misleadingly) to some
  researchers that response times ought to be collected in discrete
  time, rather than continuous time; doing so would limit researchers'
  ability to consider alternative bin widths later during analysis. I
  suggest renaming this section ``Implications for research design and
  analysis'' and making it clear that response time data still should be
  recorded in continuous time, if multiple bin widths are to be
  considered later during analysis.
\end{enumerate}
\end{quote}

Response: We renamed the section, and make it clear that RT data should
still be recorded in continuous time, so that various bin widths can be
considered later during analysis.

\begin{quote}
\begin{enumerate}
\def\labelenumi{\arabic{enumi}.}
\setcounter{enumi}{1}
\tightlist
\item
  Related to the above comment on bin width, I also think some
  discussion of the sensitivity of results and conclusions to the choice
  of bin width is both necessary and useful, given that this is a
  tutorial paper and choosing the width of the bin is a key step in
  modeling. The authors state generally that the ``goal is to find the
  smallest bin width that is supported by the amount of data available''
  . This guidance may be difficult to translate to actionable steps,
  particularly by a researcher who is new to EHA. For example, what sort
  of metric is used to determine ``support''? Is it noise in the
  estimated hazard function or another criterion? More specificity here
  could be useful, and I would recommend including a supplementary
  example that illustrates both choosing a bin width and any potential
  impact of that bin width on the conclusions.
\end{enumerate}
\end{quote}

Response: We include a supplementary example that illustrates how
results are sensitive to the choice of bin width.

\begin{quote}
\begin{enumerate}
\def\labelenumi{\arabic{enumi}.}
\setcounter{enumi}{2}
\tightlist
\item
  Figure 2: In the plot of ca(t), why is there a gap between the first
  and second orange triangles? That is, why is ca(t) missing at 160ms?
\end{enumerate}
\end{quote}

Response: Because no responses were emitted in that time bin, so hazard
equals 0 and ca(t) is undefined.

\begin{quote}
\begin{enumerate}
\def\labelenumi{\arabic{enumi}.}
\setcounter{enumi}{3}
\tightlist
\item
  P. 23 line 407-409: The text states, ``First, because the first few
  bins often contain no responses, one has to select an analysis time
  window, i.e., a contiguous set of bins for which there is data for
  each participant.'' What is meant here by ``there is data for each
  participant''? And why is it important that the first few bins often
  contain no responses---isn't a lack of response at early time points
  useful data? This sentence could benefit from being re-written for
  increased clarity.
\end{enumerate}
\end{quote}

Response: We re-wrote this sentence.

\begin{quote}
\begin{enumerate}
\def\labelenumi{\arabic{enumi}.}
\setcounter{enumi}{4}
\tightlist
\item
  P. 38 line 641-643: ``We might pick a hazard ratio\ldots{[}that{]} is
  much smaller than that previously observed'' . Do you mean closer to
  1, i.e., weaker than previously observed?
\end{enumerate}
\end{quote}

Response: Yes, we mean closer to 1, and thus a hazard ration that is
weaker than previously observed.

Rich: please check this response.

\begin{quote}
Minor comments: 1. P. 26 line 446: ``bernoulli'' should be ``Bernoulli''
\end{quote}

Response: We changed this.

\begin{quote}
\begin{enumerate}
\def\labelenumi{\arabic{enumi}.}
\setcounter{enumi}{1}
\tightlist
\item
  P. 27 line 486: I would add the section reference numbers (section
  4.6.4) for Kurz 2023a (either to the main text or to the reference
  list) to help interested readers find the relevant information.
\end{enumerate}
\end{quote}

Response: We added the section reference numbers in the main text.

\begin{quote}
\begin{enumerate}
\def\labelenumi{\arabic{enumi}.}
\setcounter{enumi}{2}
\tightlist
\item
  P. 27 line 489: If you first round a number to `digits = 1' and then
  format using `nsmall = 2', you may get a misleading value of 0 (i.e.,
  0.009 would appear as 0.00, when in fact 0.01 is more correct). This
  should be corrected.
\end{enumerate}
\end{quote}

Response: We corrected this.

\begin{quote}
\begin{enumerate}
\def\labelenumi{\arabic{enumi}.}
\setcounter{enumi}{3}
\tightlist
\item
  P. 32 line 558: ``R package lme4()'' should be ``R package lme4'' ,
  otherwise the reference appears to be for a function, rather than an R
  package. Same comment for ``faux\{\}''---``faux'' would be clearer.
\end{enumerate}
\end{quote}

Response: We followed both suggestions.

\section{Reviewer 3}\label{reviewer-3}

\begin{quote}
Comments to the Author This tutorial provides guidance on wrangling
cognitive reaction time data and addressing questions about temporal
dynamics of cognition using event history analysis (EHA) and
speed/accuracy tradeoff (SAT) analysis, with illustrations of hazard and
conditional accuracy functions. The manuscript is improved in some
respects (e.g., shortening the text by moving the frequentist analyses
to the supplement).
\end{quote}

Response: We want to thank this reviewer for his/her supportive
comments.

\begin{quote}
At the same time, the tutorial remains challenging to follow,
particularly for readers less familiar with EHA and SAT. To enhance
accessibility, I encourage the authors to ensure that key information
such as descriptions of research questions, RT data, and analytic
approaches (including equations) is presented in the main text rather
than moved to supplemental materials. This could likely be achieved
without increasing the word count by reducing redundancies and
tightening language in several places. Overall, the tutorial covers
important concepts, but there is still room to strengthen its
contributions, accessibility, and clarity.
\end{quote}

Response. We now present the descriptions of research questions, RT
data, and analytic approaches (including equations) in the main text.

\begin{quote}
Contribution of this tutorial: Thank you for clarifying the aims of the
paper (line 89). To strengthen this section, it would be helpful to
provide more rationale on the unique characteristics of RT data that
necessitate a dedicated EHA tutorial. In particular, please address
directly why a tutorial focused on RT data is needed. At present, key RT
data characteristics (including the definition of ``trial'') appear only
in the supplement. Given the centrality of these features, I recommend
moving them into the main text so the tutorial stands on its own.
\end{quote}

Response: We explain why a tutorial focused on RT data is needed, and
define key concepts (such as trial) in the main text now.

\begin{quote}
The discussion at line 93 would benefit from greater emphasis on why a
researcher might conduct EHA analyses. More clearly outlining the types
of research questions that survival analysis of RT data can address
would make the tutorial more useful and accessible. While some research
questions are introduced later (pg. 10), I recommend stating them more
clearly and incorporating them earlier (e.g., in Section 2.1).
\end{quote}

Response: We now state research questions more clearly and earlier.

\section{response from Claude:}\label{response-from-claude}

\subsection{Research Questions Addressable with Event History Analysis
in RT
Research}\label{research-questions-addressable-with-event-history-analysis-in-rt-research}

Event history analysis (survival analysis) for response times offers
unique advantages over traditional RT analysis (mean/median
comparisons). Here are key research questions it can address:

\begin{center}\rule{0.5\linewidth}{0.5pt}\end{center}

\subsection{\texorpdfstring{\textbf{1. Distributional
Questions}}{1. Distributional Questions}}\label{distributional-questions}

\textbf{How does the entire RT distribution change, not just central
tendency?}

Traditional analysis focuses on means/medians, but event history
analysis examines: - \textbf{Early vs.~late portions of the
distribution} - Do experimental manipulations affect fast vs.~slow
responses differently? - \textbf{Tail behavior} - Does a manipulation
increase probability of very slow responses? - \textbf{Skewness changes}
- Does a condition shift the entire distribution or just stretch the
tail?

\textbf{Example:} Does cognitive load slow all responses uniformly, or
does it primarily increase the probability of very slow responses while
leaving fast responses unaffected?

\begin{Shaded}
\begin{Highlighting}[]
\CommentTok{\# Examine time{-}varying hazard ratios}
\CommentTok{\# If HR \textgreater{} 1 at early times but HR \textless{} 1 at late times,}
\CommentTok{\# effect differs across the distribution}
\end{Highlighting}
\end{Shaded}

\begin{center}\rule{0.5\linewidth}{0.5pt}\end{center}

\subsection{\texorpdfstring{\textbf{2. Time-Dependent
Effects}}{2. Time-Dependent Effects}}\label{time-dependent-effects}

\subsubsection{\texorpdfstring{\textbf{When during processing does an
experimental factor exert its
influence?}}{When during processing does an experimental factor exert its influence?}}\label{when-during-processing-does-an-experimental-factor-exert-its-influence}

Hazard models can reveal \textbf{timing} of effects: - \textbf{Early
effects} (high hazard early) suggest perceptual or automatic processes -
\textbf{Late effects} suggest deliberative, controlled processing -
\textbf{Crossing hazards} indicate time-dependent or strategic effects

\textbf{Example:} Does semantic priming affect early automatic retrieval
(increased hazard at 300-500ms) or later strategic processing (decreased
hazard at 700-1000ms)?

\begin{center}\rule{0.5\linewidth}{0.5pt}\end{center}

\subsection{\texorpdfstring{\textbf{3. Speed-Accuracy
Tradeoffs}}{3. Speed-Accuracy Tradeoffs}}\label{speed-accuracy-tradeoffs}

\subsubsection{\texorpdfstring{\textbf{How do manipulations affect the
relationship between response speed and
accuracy?}}{How do manipulations affect the relationship between response speed and accuracy?}}\label{how-do-manipulations-affect-the-relationship-between-response-speed-and-accuracy}

By modeling correct and error responses separately (competing risks): -
\textbf{Conditional accuracy functions} - Accuracy as a function of RT -
\textbf{Divergence points} - When do accuracy differences emerge? -
\textbf{Strategic shifts} - Do participants adjust response caution?

\textbf{Example:} Does time pressure increase hazard equally for correct
and error responses (pure speed effect), or does it disproportionately
increase error hazard (speed-accuracy tradeoff)?

\begin{Shaded}
\begin{Highlighting}[]
\CommentTok{\# Competing risks model}
\NormalTok{model\_correct }\OtherTok{\textless{}{-}} \FunctionTok{coxph}\NormalTok{(}\FunctionTok{Surv}\NormalTok{(RT, correct) }\SpecialCharTok{\textasciitilde{}}\NormalTok{ condition }\SpecialCharTok{+}\NormalTok{ (}\DecValTok{1}\SpecialCharTok{|}\NormalTok{subject))}
\NormalTok{model\_error }\OtherTok{\textless{}{-}} \FunctionTok{coxph}\NormalTok{(}\FunctionTok{Surv}\NormalTok{(RT, }\SpecialCharTok{!}\NormalTok{correct) }\SpecialCharTok{\textasciitilde{}}\NormalTok{ condition }\SpecialCharTok{+}\NormalTok{ (}\DecValTok{1}\SpecialCharTok{|}\NormalTok{subject))}
\end{Highlighting}
\end{Shaded}

\begin{center}\rule{0.5\linewidth}{0.5pt}\end{center}

\subsection{\texorpdfstring{\textbf{4. Individual Differences in
Processing
Dynamics}}{4. Individual Differences in Processing Dynamics}}\label{individual-differences-in-processing-dynamics}

\subsubsection{\texorpdfstring{\textbf{Do individuals differ in how
effects unfold over
time?}}{Do individuals differ in how effects unfold over time?}}\label{do-individuals-differ-in-how-effects-unfold-over-time}

Random effects on time-varying components: - \textbf{Individual hazard
shapes} - Does baseline processing speed vary? - \textbf{Individual
effect timing} - Does a manipulation affect people at different
processing stages? - \textbf{Strategy differences} - Do some people show
early effects while others show late effects?

\textbf{Example:} In Stroop tasks, do high vs.~low working memory
individuals show interference at different time points?

\begin{center}\rule{0.5\linewidth}{0.5pt}\end{center}

\subsection{\texorpdfstring{\textbf{5. Sequential and Context
Effects}}{5. Sequential and Context Effects}}\label{sequential-and-context-effects}

\subsubsection{\texorpdfstring{\textbf{How do previous trials influence
current response
timing?}}{How do previous trials influence current response timing?}}\label{how-do-previous-trials-influence-current-response-timing}

Time-varying covariates capture dynamic adjustments: -
\textbf{Post-error slowing} - How long does slowing persist after
errors? - \textbf{Conflict adaptation} - Does previous trial congruency
affect current hazard shape? - \textbf{Priming dynamics} - Time course
of repetition priming across lags

\textbf{Example:} After an incongruent trial, does the hazard for the
next trial decrease uniformly, or is there increased probability of very
cautious (slow) responses?

\begin{center}\rule{0.5\linewidth}{0.5pt}\end{center}

\subsection{\texorpdfstring{\textbf{6. Censoring and Truncation
Issues}}{6. Censoring and Truncation Issues}}\label{censoring-and-truncation-issues}

\subsubsection{\texorpdfstring{\textbf{How to handle response deadlines,
anticipations, or missing
responses?}}{How to handle response deadlines, anticipations, or missing responses?}}\label{how-to-handle-response-deadlines-anticipations-or-missing-responses}

Event history analysis naturally handles: - \textbf{Response deadlines}
- Trials with no response before cutoff (right-censored) -
\textbf{Anticipatory responses} - Extremely fast responses
(left-truncation) - \textbf{Task switching} - Trials where participant
abandons task

\textbf{Example:} In adaptive testing, how do difficulty adjustments
affect the hazard of responding before the 2-second deadline?

\begin{center}\rule{0.5\linewidth}{0.5pt}\end{center}

\subsection{\texorpdfstring{\textbf{7. Cumulative Processing
Effects}}{7. Cumulative Processing Effects}}\label{cumulative-processing-effects}

\subsubsection{\texorpdfstring{\textbf{What is the probability of
responding by time t under different
conditions?}}{What is the probability of responding by time t under different conditions?}}\label{what-is-the-probability-of-responding-by-time-t-under-different-conditions}

Survival functions answer: - \textbf{Threshold timing} - When have 50\%
of responses occurred? - \textbf{Cumulative advantage} - How much ``head
start'' does one condition have? - \textbf{Processing deadlines} -
Probability of completing processing by a fixed time

\textbf{Example:} By 500ms, what proportion of participants have
completed lexical access in high vs.~low frequency words?

\begin{center}\rule{0.5\linewidth}{0.5pt}\end{center}

\subsection{\texorpdfstring{\textbf{8. Non-Proportional
Effects}}{8. Non-Proportional Effects}}\label{non-proportional-effects}

\subsubsection{\texorpdfstring{\textbf{Do effects change in magnitude or
direction over
time?}}{Do effects change in magnitude or direction over time?}}\label{do-effects-change-in-magnitude-or-direction-over-time}

Unlike ANOVA, can detect when: - \textbf{Effects emerge gradually} -
Hazard ratio increases over time - \textbf{Effects reverse} - Early
disadvantage becomes late advantage - \textbf{Effects saturate} - Strong
early effect diminishes later

\textbf{Example:} Does emotional content produce an early attention
capture (high hazard at 200-400ms) followed by later avoidance (low
hazard at 600-1000ms)?

\begin{Shaded}
\begin{Highlighting}[]
\CommentTok{\# Test time × condition interaction}
\NormalTok{model }\OtherTok{\textless{}{-}} \FunctionTok{coxph}\NormalTok{(}\FunctionTok{Surv}\NormalTok{(RT, event) }\SpecialCharTok{\textasciitilde{}}\NormalTok{ condition }\SpecialCharTok{*} \FunctionTok{log}\NormalTok{(RT) }\SpecialCharTok{+} \FunctionTok{cluster}\NormalTok{(subject))}
\end{Highlighting}
\end{Shaded}

\begin{center}\rule{0.5\linewidth}{0.5pt}\end{center}

\subsection{\texorpdfstring{\textbf{9. Mixture/Multiple Process
Questions}}{9. Mixture/Multiple Process Questions}}\label{mixturemultiple-process-questions}

\subsubsection{\texorpdfstring{\textbf{Are there distinct processing
routes or
strategies?}}{Are there distinct processing routes or strategies?}}\label{are-there-distinct-processing-routes-or-strategies}

Mixture models identify: - \textbf{Fast guesses vs.~systematic
responses} - Bimodal hazard functions - \textbf{Direct vs.~mediated
retrieval} - Multiple hazard peaks - \textbf{Parallel race models} -
Which process ``wins''?

\textbf{Example:} In recognition memory, separate hazard functions for
recollection-based (late) vs.~familiarity-based (early) responses.

\begin{center}\rule{0.5\linewidth}{0.5pt}\end{center}

\subsection{\texorpdfstring{\textbf{10. Repeated Measures
Dynamics}}{10. Repeated Measures Dynamics}}\label{repeated-measures-dynamics}

\subsubsection{\texorpdfstring{\textbf{How do within-subject effects
evolve across the
experiment?}}{How do within-subject effects evolve across the experiment?}}\label{how-do-within-subject-effects-evolve-across-the-experiment}

Multilevel event history analysis for: - \textbf{Learning/practice
effects} - Changes in hazard shape across blocks - \textbf{Fatigue
effects} - Time-on-task influences - \textbf{Within-subject stability} -
Do individual hazard patterns persist?

\textbf{Example:} Does the congruency effect (in hazard terms) decrease
across experimental blocks as participants learn to ignore distractors?

\begin{center}\rule{0.5\linewidth}{0.5pt}\end{center}

\subsection{\texorpdfstring{\textbf{11. Boundary Conditions of
Effects}}{11. Boundary Conditions of Effects}}\label{boundary-conditions-of-effects}

\subsubsection{\texorpdfstring{\textbf{Under what temporal constraints
do effects appear or
disappear?}}{Under what temporal constraints do effects appear or disappear?}}\label{under-what-temporal-constraints-do-effects-appear-or-disappear}

Examine effect boundaries: - \textbf{Minimum processing time} - When can
an effect first be detected? - \textbf{Maximum duration} - When do
effects no longer manifest? - \textbf{Critical windows} - Temporal
windows where effects are strongest

\textbf{Example:} Masked priming effects may only appear in responses
between 400-600ms, suggesting automatic but not strategic influence.

\begin{center}\rule{0.5\linewidth}{0.5pt}\end{center}

\subsection{\texorpdfstring{\textbf{Concrete Research
Examples}}{Concrete Research Examples}}\label{concrete-research-examples}

\subsubsection{\texorpdfstring{\textbf{Lexical
Decision:}}{Lexical Decision:}}\label{lexical-decision}

\emph{Does word frequency affect early perceptual processing or later
decision processes?} - High frequency may increase hazard uniformly
(proportional hazard) - Or may specifically reduce probability of very
slow responses (non-proportional)

\subsubsection{\texorpdfstring{\textbf{Flanker/Simon
Tasks:}}{Flanker/Simon Tasks:}}\label{flankersimon-tasks}

\emph{When does interference from incompatible stimuli emerge?} -
Compatible and incompatible trials may have similar early hazards
(automatic capture) - Hazards diverge at 350ms (when conflict is
detected)

\subsubsection{\texorpdfstring{\textbf{Go/No-Go:}}{Go/No-Go:}}\label{gono-go}

\emph{How does inhibition develop over time?} - Initially high hazard
for both Go and No-Go - No-Go hazard drops sharply after 200ms
(inhibition)

\subsubsection{\texorpdfstring{\textbf{Recognition
Memory:}}{Recognition Memory:}}\label{recognition-memory}

\emph{Do ``remember'' vs.~``know'' judgments have different temporal
signatures?} - ``Know'' responses: earlier hazard peak
(familiarity-fast) - ``Remember'' responses: later, broader hazard
(recollection-slow)

\begin{center}\rule{0.5\linewidth}{0.5pt}\end{center}

\subsection{\texorpdfstring{\textbf{Why Not Just Use Traditional RT
Analysis?}}{Why Not Just Use Traditional RT Analysis?}}\label{why-not-just-use-traditional-rt-analysis}

Traditional methods (t-tests on means, ANOVA) answer: \emph{``Is there a
difference in average RT?''}

Event history analysis additionally answers: - \emph{``When does the
difference emerge?''} - \emph{``Does it affect all responses or just
slow ones?''} - \emph{``How does the full distribution change?''} -
\emph{``Are there multiple underlying processes?''}

This makes it particularly powerful for \textbf{theory testing} about
cognitive mechanisms and their time course, going beyond simple
detection of effects to understanding the \textbf{dynamics} of
information processing.

\subsection{end of claude's response.}\label{end-of-claudes-response.}

\begin{quote}
In addition, the response to my previous comment that research questions
were not expanded because they are ``highly dependent on the field of
study'' seems at odds with the stated goal of offering a tutorial
specific to RT data. Providing a focused set of examples would help
balance breadth with depth and strengthen the tutorial's clarity and
relevance.
\end{quote}

Response: We discuss research questions in the main text.

\begin{itemize}
\item
  shape hazard function
\item
  differences between conditions
\item
  effects of longer time scales
\end{itemize}

\begin{quote}
Use of discrete-time EHA: Pg. 8 -- The rationale given (``as a lab, and
mainly for practical reasons, we have much more experience using
discrete-time EHA'') is not a compelling justification for the analytic
choice and feels out of place in a manuscript. One option to consider is
to revise this paragraph to instead (1) clarify the distinction between
discrete- and continuous-time approaches, (2) note that there are no
strict cutoffs for when data should be treated as discrete versus
continuous, and (3) highlight the relative advantages, limitations, or
tradeoffs of each. This framing would be more informative for readers
and align better with the tutorial's didactic purpose.
\end{quote}

Response: We revised this paragraph to (1) clarify the distinction
between discrete- and continuous-time approaches, (2) note that there
are no strict cutoffs for when data should be treated as discrete versus
continuous, and (3) highlight the relative advantages, limitations, or
tradeoff of each.

\begin{quote}
Clarity and accuracy: Pg. 3 -- The speed/accuracy tradeoff analysis is
mentioned without explanation. Even if this concept is familiar to many
experimental psychologists (as noted in your prior responses), it would
strengthen the manuscript to include one or two sentences defining key
terms or analysis types when first introduced. This would improve
accessibility for a broader readership and support the tutorial's
clarity. I highlight a few examples below, but this should be considered
throughout.
\end{quote}

Response: We explain SAT analysis early in the revised ms.

\begin{quote}
Pg. 8 -- The phrase ``discrete or continuous time units'' is introduced
without explanation. Please add a brief description of the distinction
between these two approaches.
\end{quote}

Response: We discuss the distinction now on page x of the revised ms.

\begin{quote}
Pgs. 8-10 -- Terms such as ``hazard of response'' (pg. 8) and
``discrete-time hazard function'' (pg. 9) are introduced without
definition.
\end{quote}

Response: We introducte definitions for these terms.

\begin{quote}
Likewise, Figure 1 presents hazard and conditional accuracy equations
but does not walk readers through their meaning or interpretation. I
encourage you to keep the equations in the main text and provide
accessible explanations, as this will better equip readers who may be
less comfortable with mathematical notation.
\end{quote}

Response: We keep the equations for hazard and conditional accuracy in
the main text and provide accessible explanations.

\begin{quote}
Similarly it would be helpful to provide key definitions of ``hazard''
and ``conditional accuracy estimates'' should also be provided up front,
rather than deferred to Section 2.1.2.
\end{quote}

Response: We profide the definitions on page x of the revised ms.

\begin{quote}
Pg 9 -- in discrete time, there is not instantaneous risk. Please
revise.
\end{quote}

Response: We removed ``instantaneous'' when talking about risk now.

\begin{quote}
Pg 14 -- The manuscript would be strengthened if key EHA elements were
defined directly in the main text rather than directing readers to the
supplemental material.
\end{quote}

Response: We define key EHA elements directly in the main text now.

\begin{quote}
Pg 15 -- In my original review I requested that right and left censoring
be more clearly defined. Right censoring is still referred to several
times in the manuscript without initial explanation of what this is. For
the benefit of novice readers (since this is an AMPPS tutorial), I
recommend including these definitions in the main text rather than the
supplement.
\end{quote}

Response: We include definitions for right- and left-censoring in the
main text.

\end{document}
